\chapter{Installation}
\label{ch:install}

\section{Requirements}
\begin{itemize}
  \item \textbf{OpenFOAM v2506} from openfoam.com (the ESI/OpenCFD version). The
        openfoam.org versions (OpenFOAM-11, -12, \ldots) are \emph{not} supported.
  \item C, C++ and Fortran compilers, \kw{make}, CMake and git.
  \item FFTW3 (the static library \file{libfftw3.a} must be present) and GSL.
  \item Linux: Ubuntu or Debian (also Windows Subsystem for Linux), RHEL, Rocky or Fedora,
        or an HPC cluster with environment modules.
\end{itemize}
\wtwof{}, Grid2Grid, FloatStepper and MoorDyn are part of the repository or are downloaded
by the build script; they need no separate installation. HOS-NWT or HOS-Ocean are needed
only to generate new HOS wave fields (Chapter~\ref{ch:hos}).

\section{Installing}
\begin{lstlisting}[style=shell,caption={Installing FoamWaveTank.},label={lst:install}]
# 1. Load OpenFOAM v2506 in the terminal
source /usr/lib/openfoam/openfoam2506/etc/bashrc   # Ubuntu/WSL package
# module load openfoam.com/2506                     # HPC cluster (names differ)

# 2. Get the code
git clone https://github.com/Saliyar/FoamWaveTankv2.git
cd FoamWaveTankv2

# 3. Compile (10-20 minutes)
./Allwmake
\end{lstlisting}

The solvers and libraries are installed in \file{\$FOAM\_USER\_APPBIN} and
\file{\$FOAM\_USER\_LIBBIN}, so they are found by OpenFOAM without further settings.

\section{What \texttt{./Allwmake} does}
\subsection{Prerequisite check}
Before compiling anything, \file{Allwmake} determines the platform:
\begin{description}
  \item[server] an HPC cluster with environment modules (the \kw{module} command exists);
  \item[ubuntu] Ubuntu or Debian, including WSL;
  \item[rhel] RHEL, Rocky, CentOS or Fedora;
  \item[other] anything else.
\end{description}
It then checks for the compilers, \kw{make}, CMake, git, FFTW (\file{libfftw3.a}), GSL (by
compiling and linking a small test program) and a loaded OpenFOAM v2506. On a cluster it
first tries to load an FFTW module itself.

If something is missing, \emph{nothing is compiled}. The script lists what is missing and
prints the command that installs it for the detected platform. On Ubuntu it also offers to run
that command:
\begin{lstlisting}[style=shell,caption={Typical output on a fresh Ubuntu installation.}]
Platform: ubuntu  (ubuntu)
MISSING: gcc g++ make gfortran cmake git FFTW(libfftw3.a) GSL OpenFOAM(not loaded)

  sudo apt update && sudo apt install -y build-essential gfortran cmake git \
       libfftw3-dev libgsl-dev
Install these packages now with sudo apt? [y/N]
\end{lstlisting}
If OpenFOAM is installed but not loaded, the script names the \file{etc/bashrc} file to
source. If it is not installed at all, on Ubuntu it prints the commands that install
OpenFOAM v2506 from the openfoam.com package repository:
\begin{lstlisting}[style=shell,caption={Installing OpenFOAM v2506 on Ubuntu or WSL.}]
curl -s https://dl.openfoam.com/add-debian-repo.sh | sudo bash
sudo apt install -y openfoam2506-default
source /usr/lib/openfoam/openfoam2506/etc/bashrc
\end{lstlisting}
Install what is listed and run \file{./Allwmake} again.

\subsection{Build}
When all prerequisites are found, \file{Allwmake}
\begin{enumerate}
  \item creates \file{bin/bashrc} from \file{bin/bashrc.org} if it does not exist and sources
        it; \kw{WAVES\_DIR} is set to the repository folder automatically;
  \item fetches the FloatStepper submodule;
  \item builds Grid2Grid with CMake and FFTW (\file{ThirdParty/});
  \item builds the \wtwof{} libraries, including the HOS and current theories;
  \item builds the solvers and utilities;
  \item builds FloatStepper, MoorDyn v1 and v2, and the MoorDyn restraints for
        \kw{sixDoFRigidBodyMotion}.
\end{enumerate}

\subsection{Result check}
At the end \file{Allwmake} checks that the solvers (\app{waveFoam}, \app{waveIsoFoam},
\app{floatStepper}, \app{floatStepperWaveIsoFoam}), the main utilities and the libraries
(\file{libwaves2Foam}, \file{libGrid2Grid}, \file{libmoordyn\_v1}, \file{libmoordyn},
\file{libmoorDynR1SixDoF}, \file{libmoorDynR2SixDoF}) exist. It ends with either
\begin{lstlisting}[style=shell]
COMPILATION DONE
    All solvers and libraries are in <FOAM_USER_APPBIN> and <FOAM_USER_LIBBIN>
\end{lstlisting}
or with \kw{COMPILATION FAILED} and the list of parts that were not built. In that case,
search the build log for the first line containing \kw{error}. The exit status of
\file{Allwmake} is non-zero when the build failed, so it can be used in scripts.

\section{Updating}
Bug fixes are pushed to the \kw{master} branch. To update:
\begin{lstlisting}[style=shell,caption={Updating an existing installation.}]
cd FoamWaveTankv2
git pull
./Allwmake
\end{lstlisting}
\file{Allwmake} also updates FloatStepper to its latest version. \file{./Allwclean} removes all
build files if a completely fresh build is needed.

\section{Troubleshooting}
\begin{description}
  \item[``OpenFOAM v2506 is not loaded''] Source the OpenFOAM \file{etc/bashrc} in the same
    terminal before \file{./Allwmake} (Code Frame~\ref{lst:install}).
  \item[Wrong OpenFOAM flavour] If the script reports an openfoam.org version, install
    OpenFOAM v2506 from openfoam.com.
  \item[FFTW not found] Install the development package (\kw{libfftw3-dev} on Ubuntu,
    \kw{fftw-devel} and \kw{fftw-static} on RHEL) or load an FFTW module on a cluster. A
    custom FFTW can be used by setting \kw{USER\_FFTW3\_LIB\_PATH="<folder with libfftw3.a>"}
    near the top of \file{Allwmake}.
  \item[Solver not found after building] Open a new terminal, source OpenFOAM, and check that
    \kw{\$FOAM\_USER\_APPBIN} is the folder reported by \file{Allwmake}.
\end{description}
